% Generated draft: requires mathematical and source review.
\documentclass[11pt]{article}
\usepackage[margin=25mm]{geometry}
\usepackage{fontspec}
\setmainfont{Latin Modern Roman}
\usepackage{amsmath,amssymb,mathtools,unicode-math}
\setmathfont{Latin Modern Math}
\usepackage{enumitem,needspace,xurl,hyperref}
\setlength{\parindent}{0pt}
\setlength{\parskip}{5pt}
\setlength{\emergencystretch}{3em}
\newcommand{\N}{\mathbb{N}}
\newcommand{\Z}{\mathbb{Z}}
\newcommand{\Q}{\mathbb{Q}}
\newcommand{\R}{\mathbb{R}}
\newcommand{\C}{\mathbb{C}}
\newcommand{\F}{\mathbb{F}}
\newcommand{\A}{\mathbb{A}}
\newcommand{\PP}{\mathbb{P}}
\newcommand{\E}{\mathbb{E}}
\newcommand{\eps}{\varepsilon}
\newcommand{\dd}{\,\mathrm{d}}
\begin{document}
\section*{UnsolvedMath ID: 20002484}
\textbf{Source problem code:} catalogue entry
\subsection{Definitions and mathematical statement}
Let \((M^n,g)\) be a compact smooth Riemannian manifold without boundary. A real \(q\)-current is a continuous linear functional on smooth \(q\)-forms. It is closed if \(T(d\eta)=0\) for every smooth \((q-1)\)-form \(\eta\); closed currents, rather than arbitrary currents, determine homology classes. Use pointwise comass
\[
 \|\omega_x\|_{\mathrm{comass}}=
 \sup\{|\omega_x(v_1,\ldots,v_q)|:v_1,\ldots,v_q\text{ orthonormal}\},
\]
and mass \(\mathbf M(T)=\sup\{T(\omega):\sup_x\|\omega_x\|_{\mathrm{comass}}\le1\}\). A closed finite-mass current is normal. Minimizing mass among real normal representatives of \(h\in H_q(M;\R)\) defines its stable norm. Minimizing among integral rectifiable currents is a different admissible class and requires integral homology data.

Victor Bangert's question asks for the structure of homologically mass-minimizing currents in intermediate degrees \(1<q<n-1\). The source contrasts geodesic descriptions in degree one and measured minimizing hypersurface laminations in codimension one, possibly with singularities. Its intermediate-dimensional question is open ended, with no one proposed regularity theorem or explicit coefficient category. The comass convention and real-versus-integral distinction are essential.

The aid's flat-torus result works on the unit-volume \(T^n=\R^n/\Z^n\). For a constant \(q\)-vector \(v\), define
\[
 T_v(\omega)=\int_{T^n}\omega_x(v)\,\dd x.
\]
It is a closed real normal current. If \(|v|_{\mathrm{mass}}\) denotes the algebraic mass norm dual to comass, the aid asserts
\[
 \mathbf M(T_v)=|v|_{\mathrm{mass}}=\|[T_v]\|_{\mathrm{stable}}.
\]
A constant closed calibrating form of comass one attaining the dual norm proves minimality; the norm is not silently replaced by the Euclidean norm on all \(q\)-vectors.

For every \(n\ge4\) and \(2\le q\le n-2\), let
\[
 v=e_1\wedge\cdots\wedge e_{q-2}\wedge
 (e_{q-1}\wedge e_q+e_{q+1}\wedge e_{q+2}),
\]
where the initial wedge is one when \(q=2\). The constant form
\[
 \theta=dx_1\wedge\cdots\wedge dx_{q-2}\wedge
 (dx_{q-1}\wedge dx_q+dx_{q+1}\wedge dx_{q+2})
\]
has comass one and \(\theta(v)=2\). Thus \(T_v\) has mass two and minimizes its integral homology class viewed over \(\R\). Its mass distribution is diffuse with respect to \(n\)-dimensional volume, so it is not a rectifiable \(q\)-current. More strongly, its nondecomposable polar \(q\)-vector prevents it from being the current of one oriented measured \(q\)-lamination with one tangent \(q\)-plane almost everywhere. This does not forbid a sum of currents from different laminations.

The same integral class has an integral calibrated minimizer: the sum of the coordinate \(q\)-tori in directions \((e_1,\ldots,e_q)\) and \((e_1,\ldots,e_{q-2},e_{q+1},e_{q+2})\), each with positive orientation. Their intersection has lower dimension, both are calibrated by \(\theta\), and the sum has mass two. Thus a class admitting an integral minimizer can simultaneously admit a diffuse nonlaminar real minimizer. The obstruction is to extending the endpoint lamination picture to \emph{all real minimizers}, not a general impossibility of rectifiable or integral minimizing representatives.
\subsection{Short English statement}
The intermediate-dimensional structure depends on the current category and the comass mass norm. The aid gives, in every intermediate degree, a flat-torus integral class with both a mass-two integral minimizer and a mass-two diffuse real minimizer with nondecomposable tangent polar, obstructing a universal lamination description.
\subsection{Sources}
\begin{itemize}
\item Local registry, permanent UnsolvedMath ID \(20002484\), exact public code \texttt{catalogue entry}; canonical location \texttt{aim-pdes-notes.json}, \texttt{notes[95]}.
\item Victor Bangert, Problem 27 in AIM \emph{New connections between dynamical systems and PDE's}, notes by S. Bolotin; supplied reference \url{https://aimath.org/WWN/dynpde/dynpde.pdf}. The verification field reports that OCR lost the initial 2, producing source item 7, and that original AIM HTML confirms the question. Neither source was accessed here.
\item Supplied research attempt 1: \texttt{partial\_result}, run and registry status \texttt{partially\_solved}, rigor label \texttt{proved}, low novelty confidence; classification \texttt{PARTIAL-PROGRESS}, created and updated 14 August 2026. No proposer year, original remarks, literature entry, or primary citation for the flat-torus construction is supplied.
\item The aid explicitly states that it is machine generated, not peer reviewed, and requires independent verification before citation.
\end{itemize}
\end{document}
